\clearpage
\section{项目详情与使用方式}

\subsection{问题定义：转写不是学习过程的终点}

纸面资料的主要问题不是缺少保存手段，而是保存之后难以形成后续动作。一张板书照片进入相册后，通常只剩下拍摄时间和图片内容两个线索。复习者需要重新阅读整页，手工判断其中包含哪些知识点，再寻找相近题目。答错以后，还要自行定位是概念、变形还是计算出了问题。通用 OCR 与公式识别服务改善了转写质量，但它们交付的对象仍是文本或 LaTeX，并不管理练习、作答与错误诊断。

本作品将输入视为学习资料的开始，而不是识别任务的结束。识别结果先成为带来源的卡片，卡片再成为命题和检索的共同材料。题目保留对应卡片编号，作答写入尝试记录，追问增加卡片热度。这样，同一份资料能够在不同时间被继续使用，不需要在多个工具之间复制文本。系统的价值主要由这条数据路径决定，而不是由某个模型的单次回答决定。

这一定位也限定了技术范围。作品不自行实现通用 OCR 引擎，不训练或微调视觉、语言基础模型，不将免费样例演示包装成真实识别能力。可替换的 Provider 负责模型调用，本地代码负责数据结构、验算、判分、检索、事件和交互。对于无法机械验证的开放题，系统只给出有限检查并明确标记，避免把语言模型的自信当作正确性证明。

\subsection{目标用户与可验收需求}

目标场景来自开发者自身的学习过程与对现有工具的试用，没有开展大样本用户研究。因此，以下场景用于说明设计假设，不作为学习效果已经得到验证的证据。当前版本优先服务个人数学资料整理，教师出题属于同一能力的复用，不包含班级管理、学生账号和权限分发。

\begin{table}[H]
\centering\footnotesize
\caption{目标场景、现有断点与本作品的验收方式}
\begin{tabularx}{\textwidth}{@{}L{2.4cm}L{4.3cm}Y L{4.0cm}@{}}
\toprule
\textbf{使用者} & \textbf{典型任务} & \textbf{原有断点} & \textbf{当前可验收结果} \\
\midrule
理工科学生 & 课后整理一页函数板书并练习 & 图片难检索，原题缺少变式，错因靠自查 & 至少生成一张卡片；得到带验算级别的练习；答错后出现讲解 \\
高中生 & 复习错题并追问变形依据 & 手抄耗时，通用回答缺少本题出处 & 卡片可编辑；追问返回引用卡片；修正历史可回看 \\
教师或助教 & 从一份讲义生成随堂练习 & 改数字后仍需逐题核算答案 & 一次生成 1--12 题；不合格题被驳回；验算账单与级别入库 \\
\bottomrule
\end{tabularx}
\end{table}

需求被改写为能够通过接口或测试判定的条件。上传一页资料以后，\env{GET /api/cards} 应返回结构化卡片；生成接口的 \env{count} 被限制在 1--12；题目对象必须带 \env{verify_level} 与 \env{verify_log}；作答接口返回正确性或讲解；批改接口在结束事件中交付固定字段；追问响应带 \env{citations}。无密钥、无 GPU 和无网络是缺省演示路径的约束，而不是额外功能。这个约束促使系统保留本地 Provider、样例数据和确定性规则。

\begin{table}[H]
\centering\footnotesize
\caption{功能矩阵与实现边界}
\begin{tabularx}{\textwidth}{@{}L{2.5cm}Y L{4.1cm}@{}}
\toprule
\textbf{环节} & \textbf{当前能力} & \textbf{边界与失败处理} \\
\midrule
纸面输入 & 摄像头拍照、五种图片格式；拍照最长边 1920 像素，图片上限 12 MB & 局域网纯 HTTP 下浏览器不允许摄像头时改用文件；超限返回 413 \\
文档输入 & PDF、DOCX、TXT、MD；PDF 文本层可用时直读，否则真实 VLM 下渲染前 3 页 & 文档上限 20 MB、正文 2 万字；长文档截断；mock 不渲染扫描页 \\
知识卡片 & 标题、摘要、学科、标签、公式、正文；四字段可人工修正 & 来源与创建时间不可改；不存在返回 404；同值提交不产生历史 \\
练习生成 & 自然语言描述题量、难度和题型；逐题进入练习卷 & 多给部分被裁掉，少给不编造补齐；题量限制 1--12 \\
质量检查 & sympy 代入验算、结构检查、有限重试与验算日志 & 开放题和证明题无法由当前规则证明正确；格式问题可降级并标级 \\
作答判分 & 字符串、数值、sympy 等价判定；记录作答 & 多答案、集合与区间写法仍未覆盖；解析异常退回旧规则 \\
AI 批改 & 得分、得分点、分步清单、首个错误、更优解法与评语 & mock 不做方法语义推断；模型失败时回退本地规则并标 \env{json-fallback} \\
讲解追问 & 参数化函数图；基于卡片的回答与出处；可选混合检索 & 当前只有一个滑块模板；单轮提问无对话上下文；无命中不硬凑出处 \\
知识谱系 & 共享标签建边，节点大小表示出题量，颜色表示被问次数 & 不是认知诊断模型；计数只反映本工具内的行为 \\
过程观测 & 七 Agent 看板、SSE、事件落库、按 \env{run_id} 回放 & 内存事件有容量上限；单 worker；前端不自动重连 \\
桌面交付 & 目录版与单文件版，内置 Python 和前端静态资源 & 未签名；单文件首次自解压慢；桌面包含 AGPL 组件 \\
\bottomrule
\end{tabularx}
\end{table}

\subsection{一次完整使用链路}

上传后，后端先检查后缀与大小，再用随机文件名保存。识别阶段输出标题、正文、公式和表格等块，原始结果写入 \env{source.raw_json}；萃取结果形成带来源的卡片。用户可对照原图与识别块修改标题、摘要、公式或正文，真实变化写入 \env{card_revision}。

命题 Agent 按卡片和自然语言需求生成题目与可选验算账单，质检 Agent 逐题给出结果和级别。未通过且仍有轮次时，编排器带原因返回命题，最终只交付通过当前检查的题。快速判分负责正误和答错讲解，详细批改交付得分点、步骤和首个错误；真实 Provider 的方法归因仍缺充分回归。追问先检索卡片再作答，出处单独返回，卡片计数最后进入谱系。

\begin{keybox}[为什么保留中间产物]
识别错、萃取错和命题错需要不同的处置。只保存最终文本时，用户无法判断错误从哪一步进入。保存原始块、卡片、验算日志和事件后，排障可以先核对来源，再修正卡片，最后重新生成。中间产物增加了数据库字段和接口，但减少了不可解释的重跑。
\end{keybox}

\subsection{使用方式一：Windows 桌面版}

Windows 10/11 x64 推荐运行目录版。解压交付包后双击程序，不需要预装 Python、Node.js，也不需要先申请 API Key。Electron 主进程先创建窗口，再选择空闲端口拉起内置后端，依次尝试 8710、8711、8712。健康检查就绪后才加载界面，退出时结束后端进程树。数据写入用户可写目录，配置从可执行文件同级的 \file{.env} 读取。v0.1.6 的单文件版也通过显式环境变量定位该文件，避免把自解压临时目录误认为用户目录。

程序未做代码签名，Windows SmartScreen 或安全软件可能提示风险。复现者应核对交付包哈希，再决定是否允许运行。不应以关闭安全软件作为通用解决办法。目录版与单文件版的差别主要在启动和分发，业务代码、Provider 配置和数据库结构一致。

\begin{table}[H]
\centering\footnotesize
\caption{交付形态与已测运行特征}
\begin{tabularx}{\textwidth}{@{}L{2.6cm}L{2.7cm}L{3.1cm}Y@{}}
\toprule
\textbf{形态} & \textbf{体积} & \textbf{启动实测} & \textbf{适用建议} \\
\midrule
Web 源码 & 依环境而定 & 仓库虚拟环境下健康检查约 1.2 s 就绪 & 适合审查代码、运行测试和修改功能 \\
目录版 & 压缩包约 200.5 MB，解压约 471.3 MB & 窗口 0.18--0.21 s；后端 2.5--5.1 s；界面 3.9--6.4 s & 正式演示优先；启动反馈快，配置位置稳定 \\
单文件便携版 & 约 136.8 MB & 已解压时界面约 14.5 s；首次自解压合计约 48.1 s & 适合单文件分发；首次运行需等待，不作为现场首选 \\
\bottomrule
\end{tabularx}
\end{table}

表中冷启动数字来自 v0.1.5 构建，v0.1.6 修改的是单文件版配置定位，因此不把这些数字重新宣称为 v0.1.6 全量性能复测。打包体积较大的原因是内置 Python、科学计算依赖、Electron 与前端资源。先开窗再启动后端是为了让双击操作立即有反馈，但窗口出现不等于服务已可用，界面仍以健康检查为准。

\subsection{使用方式二：Web 源码模式}

源码模式要求 Python 不低于 3.11，Node.js 不低于 18；前端单元测试使用 TypeScript 类型擦除，需要 Node.js 不低于 22.6。公开仓库地址为 \repo。取得代码后，在仓库根目录复制 \file{.env.example} 为 \file{.env}。配置文件必须位于仓库根目录，后端不会从 \file{apps/server/} 隐式读取另一份配置。

\begin{lstlisting}[language=bash,caption={零密钥最小配置（仓库根目录的 .env）},label={lst:mock-env}]
VLM_PROVIDER=mock
LLM_PROVIDER=mock
\end{lstlisting}

后端在 \file{apps/server/} 中安装 \file{requirements.txt}，测试还需要 \file{requirements-dev.txt}，随后以 \env{app.main:app} 启动 uvicorn。前端在 \file{apps/web/} 安装依赖并运行 Vite。开发地址为 \url{http://127.0.0.1:5173}，后端健康检查为 \url{http://127.0.0.1:8710/api/health}。完整命令同时列在附录 A，复现时不必从正文拼接。

启动后依次检查后端监听、健康检查、Vite 地址和七 Agent 看板，完整命令与通过标准见附录 A。自动化测试使用临时目录，不读取或清空开发数据库。

缺省 mock 并不意味着绕过业务链路。它只替换视觉和语言模型的返回值，编排、数据库、验算、判分、检索和 SSE 均照常执行。其局限也必须同时说明：视觉 mock 不读取图片内容，只按原始文件名从二次函数、一元二次方程、三角函数三套样例中选择，未识别名称时退回二次函数。文件名可以帮助复现，但不能作为识别准确率的证据。

\subsection{使用方式三：真实模型与演示重置}

当前已联调组合为智谱视觉模型和 DeepSeek 官方文本模型。真实配置只需在最小配置上增加 Provider、模型名与密钥。下面示例故意省略密钥值，实际运行时由本机 \file{.env} 或环境变量提供，仓库、日志和文档均不保存真实 Key。

\begin{lstlisting}[language=bash,caption={真实 Provider 配置示例（密钥仅留占位）},label={lst:real-env}]
VLM_PROVIDER=glm
VLM_MODEL=glm-4v-flash
VLM_API_KEY=<local-vlm-key>
LLM_PROVIDER=deepseek
LLM_MODEL=deepseek-chat
LLM_API_KEY=<local-llm-key>
\end{lstlisting}

两个 \env{*_BASE_URL} 留空时使用代码中的官方兼容端点预设。显式环境变量优先于预设，未知 Provider 不会获得虚构的模型名。健康检查会分别报告视觉和文本是否就绪，缺少 Key、地址或模型名时不会静默改成 mock。视觉免费档能证明真实链路可用，但不代表更高档位模型的识别和格式稳定性。批量生成还会受上游速率限制，现场演示应优先使用已验证样例或 mock。

真实联调中，单页识别实测约 5 s，出题第 1 轮出现 3 题驳回、2 题通过，第 2 轮全部通过。该记录用于证明驳回和重生成不是 mock 独有路径，不构成多图片、多学科统计。DeepSeek 与智谱以外的分支不在本次真实回归结论内。

演示前，源码模式可在仓库根目录运行 \file{python apps/server/scripts/reset_demo.py}。重置脚本清理卡片、题目、尝试、修订、向量缓存和事件等运行表，缺省保留上传原图，只有显式传入 \env{--purge-images} 才删除。桌面包提供“一键还原演示状态.cmd”，停止程序后清理数据库、上传文件、日志与引导标记。测试与重置使用同一组表名来源，避免新增数据表后只清掉其中一部分。

\subsection{数据位置、隐私与单用户边界}

Web 数据缺省位于仓库 \file{data/}，桌面版使用用户可写目录。原始文件按日期和标题归档，数据库与 Provider 预设均不保存密钥。mock 可完全离线；切换真实 Provider 后，任务所需图片、卡片、题目或作答会发送给所选服务，不能再宣称数据不出机。

当前系统没有登录、用户归属和班级权限，不应直接暴露公网。共享部署前必须补鉴权、隔离、配额和数据删除策略，不能把本机运行等同于多租户安全。
